\section{Statistical uncertainty and evaluation denominators}
\label{sec:uncertainty}
Table~\ref{tab:uncertainty} gives 95\% intervals for the main-text
comparisons. For a rate $p$ over $n$ tasks we use the binomial standard error
$\sqrt{p(1-p)/n}$; for a difference between two methods we combine the two
standard errors as if the rates were independent. Both methods run on the same
tasks, so a paired test would give narrower intervals; the values below are
therefore conservative. On Related-Lite99 and DeepSWE113 the rates are
five-run means; averaging runs reduces run-to-run variance but not the
variance from sampling the tasks, which the binomial bound captures.

\begin{table}[htbp]
\centering
\caption{Differences in percentage points with conservative (unpaired)
binomial 95\% half-widths. Bold marks differences that exceed their
half-width. Ablation rows are on Related-Lite99 and measure the drop from the
full system.}
\label{tab:uncertainty}
\small
\begin{tabular}{llrr}
\toprule
Benchmark & Comparison & Difference & 95\% half-width \\
\midrule
Verified500 ($n{=}500$) & ContextGraph $-$ no memory & \textbf{7.6} & 5.9 \\
 & ContextGraph $-$ ReasoningBank & 2.2 & 5.8 \\
 & Other repositories only $-$ no memory & 1.4 & 6.0 \\
 & Gemini 3.1 Pro pass@1: ContextGraph $-$ no memory & \textbf{6.6} & 4.3 \\
\midrule
Related-Lite99 ($n{=}99$) & ContextGraph $-$ no memory & \textbf{13.7} & 12.4 \\
 & ContextGraph $-$ ExpeRepair & 3.8 & 13.1 \\
 & Other repositories only $-$ no memory & 7.3 & 12.0 \\
\midrule
DeepSWE113 ($n{=}113$) & ContextGraph $-$ no memory & 11.3 & 12.2 \\
 & ContextGraph $-$ Supermemory & 5.9 & 12.0 \\
 & Other repositories only $-$ no memory & 6.7 & 12.5 \\
\midrule
Ablation (Related-Lite99) & Raw trajectories & \textbf{14.6} & 12.3 \\
 & Static playbook & \textbf{14.3} & 12.3 \\
 & Fixed top-$k$ retrieval & 10.8 & 12.6 \\
 & Static pool & 6.4 & 13.0 \\
 & No graph links & 4.0 & 13.1 \\
 & Planning-time search only & 1.9 & 13.2 \\
\bottomrule
\end{tabular}
\end{table}

\paragraph{Evaluation denominators.}
Resolution uses completed evaluations as the denominator
(Section~\ref{sec:setup}). Table~\ref{tab:denominator} reports, for each
method, how many evaluations did not complete and the resolution rate when
every task counts, with an incomplete evaluation scored as unresolved.

\begin{table}[htbp]
\centering
\caption{Incomplete evaluations and all-task resolution rates in the main
comparison. On the five-run benchmarks, counts are totals over five runs
($5\times99$ and $5\times113$ evaluations).}
\label{tab:denominator}
\small
\begin{tabular}{lcccccc}
\toprule
 & \multicolumn{2}{c}{Verified500} & \multicolumn{2}{c}{Related-Lite99} & \multicolumn{2}{c}{DeepSWE113} \\
Method & Incomplete & All-task (\%) & Incomplete & All-task (\%) & Incomplete & All-task (\%) \\
\midrule
No memory & \authorfill{} & \authorfill{} & \authorfill{} & \authorfill{} & \authorfill{} & \authorfill{} \\
Strongest baseline & \authorfill{} & \authorfill{} & \authorfill{} & \authorfill{} & \authorfill{} & \authorfill{} \\
ContextGraph & \authorfill{} & \authorfill{} & \authorfill{} & \authorfill{} & \authorfill{} & \authorfill{} \\
\bottomrule
\end{tabular}
% AUTHOR: fill from the run logs for all eight methods (or at least these
% three), and add per-run standard deviations for the five-run benchmarks.
\end{table}
